Game design space has been measured almost exclusively on a single quality axis.
We apply principal component analysis to the TIGRIS corpus (150 games, 32 axes)
and show that normalization choice is decisive: raw PCA yields PC1 = 81.3\%,
confirming a dominant general quality factor; weight normalization reduces this to 72.6\%;
within-game z-score normalization --- centering each game on its own axis mean ---
reduces PC1 to 50.0\% and requires 4 components to explain $\geq$80\% of variance.
The resulting components correspond to five theoretically coherent and empirically separable
design dimensions: Elegance/Entry (50\%), Game-Architecture (18.2\%), Tension profile (9.9\%),
Interaction geometry (6.3\%), and Range/ergodicity (residual).
These five dimensions --- named TIGER in the companion framework paper TG-B1 ---
are confirmed orthogonal by 42 independent collision votes from the adversarial
Parliament review procedure applied across the corpus.
The within-game normalization also reveals a genuine design trade-off:
Information Architecture Type (A32) negatively correlates with Count-Robustness
($r = -0.779$), confirming that hidden-identity and inverted-visibility games
belong to a structurally distinct design paradigm.
